\documentclass[aps,prl,reprint,superscriptaddress,floatfix,nofootinbib,amssymb]{revtex4-2}
%% Target: Physical Review Letters (PRL) or, on overflow, JHEP Theory-Mathematics.
%% Length budget: 4 PRL pages (~3500 words) — content compressed to dense statements +
%% sketch-proofs; full numerical evidence in companion notes
%% (see \texttt{OP\_NEGATIVE\_PAPER\_DRAFT\_2026-05-17.md}).
%%
%% Naming convention follows existing \texttt{Paper\_HSH\_v3\_letter\_JNT\_v2/main.tex}.

\usepackage{amsmath,amssymb,amsthm,mathrsfs}
\usepackage{booktabs}
\usepackage{hyperref}
\usepackage{xcolor}

\newtheorem{theorem}{Theorem}
\newtheorem{proposition}[theorem]{Proposition}
\newtheorem{lemma}[theorem]{Lemma}
\newtheorem{corollary}[theorem]{Corollary}
\theoremstyle{definition}
\newtheorem{definition}[theorem]{Definition}
\newtheorem{conjecture}[theorem]{Conjecture}
\theoremstyle{remark}
\newtheorem{remark}[theorem]{Remark}

\DeclareMathOperator{\Cl}{Cl}
\DeclareMathOperator{\rk}{rk}
\DeclareMathOperator{\Gal}{Gal}
\DeclareMathOperator{\Tr}{Tr}
\DeclareMathOperator{\sgn}{sgn}
\newcommand{\Q}{\mathbb{Q}}
\newcommand{\Z}{\mathbb{Z}}
\newcommand{\R}{\mathbb{R}}
\newcommand{\C}{\mathbb{C}}
\newcommand{\N}{\mathbb{N}}
\newcommand{\F}{\mathbb{F}}
\newcommand{\OK}{\mathcal{O}_K}
\newcommand{\msb}{\overline{\mathrm{MS}}}
\newcommand{\Tier}[1]{\textsc{(Tier #1)}}

\begin{document}

\title{No-Go Theorems for the Arithmetic Determination of the\\
Yang--Mills Mass Gap Absolute Scale}

\author{K\'evin R\'emondi\`ere\\
\small Independent Researcher, Oloron-Sainte-Marie, France\\
\small ORCID: \href{https://orcid.org/0009-0008-2443-7166}{0009-0008-2443-7166}\\
\small \texttt{kevin.remondiere@gmail.com}}
\affiliation{Independent researcher, Oloron-Sainte-Marie, France\\
ORCID: \texttt{0009-0008-2443-7166}}
\date{17 May 2026}

\begin{abstract}
We assemble four formally negative results delineating what an arithmetic
framework anchored on imaginary quadratic data \emph{cannot} predict about
the four-dimensional Yang--Mills mass gap. (T1) The absolute mass gap
$m_{\text{gap}}$ is identifiable, in the continuum limit, with the
renormalisation-group invariant scale $\Lambda_{\msb}$; under
dimensional transmutation it carries no arithmetic content. (T2) Five
candidate identifications of the string tension $\sqrt{\sigma}$ with a
classical arithmetic period (Chowla--Selberg, Borel regulator,
Lichnerowicz spectral edge, Bekenstein, Hawking--Page) are falsified at
$50$--$60$ digit \texttt{PARI/GP} precision against the
Athenodorou--Teper 2021 lattice data. (T3) The hypothetical
``topological ladder'' $D_{\mathrm{top}}(N)$ mapping each compact
$\mathrm{SU}(N)$ to a class-group anchor $2^{\rk_2 \Cl K} = N^2-1$ has
exactly one solution, $N=3$, by Mih\u{a}ilescu's theorem (the former
Catalan conjecture); no extension to $N\ge 4$ exists. (T4) The
$\sim\!10^{25}$ orders of magnitude separating the bare cosmological
energy density from $\Lambda_{\text{obs}}$ are correctly read as an
\emph{exponentiated} instanton action $\exp(-8\pi^2/g^2)$, not as a
spectral coupling gap; the two scales obey different beta-functions.
A formal Bayesian model selection between the arithmetic-determines-absolute
hypothesis ($H_0$) and the dimensionless-only hypothesis ($H_1$) yields
the aggregate $\log_{10} B_{1/0} \simeq 46$ (natural-log
$\log B \simeq 107$), with a $\pm 2\sigma$ robustness band
$\log_{10} B \in [14, 28]$ under variation of the priors; this is
equivalent to a Gaussian model-selection significance
$N_\sigma = \sqrt{2 \log B} \simeq 14.6$ in favour of the
negative interpretation, and $N_\sigma > 8$ across the entire
robustness band. We complement these with the surviving five positive
pillars (param-free $F(N)$, 165 exact cross-Galois Q-rationals,
Beilinson--Damerell pairing, HSH v3 genus formula, BSD $\varphi^3$
identity at $D=-15$) and explicit falsifiables for each negative
theorem. The theorems indicate \emph{where} the five specific candidate
identifications tested here cannot, on their own, fix the absolute
mass scale of pure Yang--Mills from class-group arithmetic; they do
not preclude all future arithmetic constructions, but they bound the
present program away from those tested routes. The remaining
constructive gap is located within the W2 (Bakry--\'Emery) wall.
\end{abstract}

\maketitle

\section{Introduction}\label{sec:intro}
The Clay Millennium problem on Yang--Mills (YM) existence
\cite{JaffeWitten2000} divides into two layers: \emph{existence} of a
mass gap $m_{\text{gap}}>0$ in 4D pure $\mathrm{SU}(N)$ continuum
theory, and \emph{quantitative} prediction of its value (in units of
the string tension $\sqrt{\sigma}$, or equivalently of
$\Lambda_{\msb}$). The lattice answer at $N=3$ is settled to several
percent: $m_{0^{++}}/\sqrt{\sigma} = 3.405(21)$
\cite{AthenodorouTeper2021}, while the absolute scale is set by the
QCD chiral input.

A recent line of work \cite{ECI_v15_2026, RemondiereHSH2026}
(``Extended Class-group Identities'', ECI v15) has accumulated
substantial empirical structure relating $\mathrm{SU}(N)$ glueball
ratios to ideal class groups of imaginary quadratic fields. The most
mature deliverables are the parameter-free closed form
$F(N) = (N^2+1)/N^2 \cdot 9/10$
matching the Athenodorou--Teper N-scan at $\chi^2/\mathrm{ndf}=1.135$
with zero free parameters; the genus-theoretic count
$r(D) = 2^{\rk_2 \Cl K}$ of Q-rational weight-$5$ CM newforms
\cite{RemondiereHSH2026}; and $165$ exact cross-Galois L-value ratios
verified to $50$--$60$ digit \texttt{PARI/GP} precision. These
results, however, control only \emph{dimensionless ratios}. The
present Letter establishes that this is no accident: four independent
no-go arguments rule out arithmetic determination of the
\emph{absolute} mass scale.

The four theorems are not deep individually --- (T1) is essentially
classical dimensional transmutation, (T2) is direct numerical
falsification, (T3) is an elementary Catalan obstruction, (T4) is a
correct reading of standard semiclassical accounting --- but jointly
they trace a sharp boundary in the program. Beyond that boundary
arithmetic predictions are precluded \emph{in principle}; within it
the program retains five clean pillars that we recap in
Sec.~\ref{sec:pillars}.

We work throughout in the standard Wilson lattice $\mathrm{SU}(N)$
setting, with bare coupling $\beta = 2N/g^2$ and 't~Hooft coupling
$\lambda = g^2 N$. All arXiv identifiers cited are live-verified
(\texttt{verify-arxiv} run 17 May 2026).

\section{Background: ECI v15 framework}\label{sec:background}
For an imaginary quadratic field $K = \Q(\sqrt{D})$ of fundamental
discriminant $D<0$, write $h_K = |\Cl(K)|$, $r = \rk_2 \Cl(K)$, $e =
\exp(\Cl(K))$, $\omega(|D|) = $ number of distinct prime divisors.
Gauss's genus theorem \cite[\S\S225--237]{Gauss1801} gives $r =
\omega(|D|) - 1$ for fundamental $D < 0$. Two theorems established in
the present program are used as inputs.

\begin{theorem}[Universal Q-rational orbit count $D'$ \cite{ECI_v15_2026}]
\label{thm:Dprime}
For odd $w \ge 3$, the number of Galois orbits of normalised CM
newforms in $S_w^{\mathrm{new}}(\Gamma_0(|D|), \chi_D)$ with
$\Q$-rational Hecke eigenvalues equals
$N_w(K) = 2^r$ if $e \mid (w-1)$, and $0$ otherwise.
\end{theorem}

\begin{theorem}[Lemma C universal \cite{ECI_v15_2026}]\label{thm:LemmaC}
The cup-product matrix on the genus subgroup of $H^2(G_{K,S}, \mu_2)
\subset \mathrm{Br}(K)[2]$ vanishes for all fundamental $D<0$ and all
$r \ge 1$ (Hasse--Brauer--Noether \cite{HBN1932} + the local
restriction theorem).
\end{theorem}

The framework's quantitative anchor on the gauge side is the
parameter-free closed form
\begin{equation}\label{eq:FN}
F(N) \;:=\; \frac{m_{0^{++}}(N)}{m_{0^{++}}(3)} \;=\; \frac{9}{10}
\cdot \frac{N^2+1}{N^2},
\end{equation}
obtained from the 't~Hooft genus expansion $F(N) =
(1+c/N^2)/(1+c/9)$ by setting the single-handle-to-planar ratio $c=1$
(see \cite{ECI_v15_2026,tHooft1974}). Equation~\eqref{eq:FN} fits
the Athenodorou--Teper $\mathrm{SU}(N)$ scan, $N\in\{2,3,4,5,6,7,8\}$
plus the large-$N$ extrapolation \cite[Tab.~34, 38]{AthenodorouTeper2021},
with $\chi^2/\mathrm{ndf} = 1.135$ ($p=0.34$), and zero free parameters
\cite{ECI_v15_2026}. Dropping $\mathrm{SU}(2)$ per
\cite[\S5.6]{AthenodorouTeper2021}'s own protocol gives
$c = 1.0007 \pm 0.063$, i.e.\ $c=1$ to four digits.

Throughout, ``$\sqrt{\sigma}$'' refers to the lattice infrared string
tension, calibrated to either Necco--Sommer ($\sqrt{\sigma}=0.444$
GeV) or AT2021 ($\sqrt{\sigma}\simeq 0.485$ GeV).

\section{Theorem 1: $m_{\text{gap}}$ is not arithmetic}\label{sec:T1}

\begin{theorem}[Absolute mass gap is non-arithmetic]\label{thm:T1}
Let $K = \Q(\sqrt D)$ be an imaginary quadratic field, with the
arithmetic data $(D, h_K, \Cl K, \rk_2, R_{\mathrm{Borel}},
\{L(\psi_\chi, s)\}_\chi)$ all assumed known to arbitrary precision.
Let $m_{\text{gap}}(N)$ denote the 4D continuum $\mathrm{SU}(N)$ pure
gauge mass gap. Then no $\R$-valued functional of the listed
arithmetic data alone equals $m_{\text{gap}}(N)$ in any physical unit
system.
\end{theorem}

\begin{proof}[Proof sketch]
By Coleman--Weinberg dimensional transmutation \cite{ColemanWeinberg1973}
and Wilson's renormalisation group \cite{Wilson1971}, the 4D
$\mathrm{SU}(N)$ effective action depends on a single dimensionless
running coupling $g(\mu)$ and on the dimensionful renormalisation
scale $\mu$. The combination
\begin{equation}
\Lambda_{\msb} \;=\; \mu \exp\!\bigl(-1/\bigl[2 b_0 g^2(\mu)\bigr]\bigr)
\, \bigl(b_0 g^2(\mu)\bigr)^{-b_1/(2 b_0^2)} (1 + O(g^2))
\label{eq:LambdaMSbar}
\end{equation}
with $b_0 = 11N/(48\pi^2)$, $b_1 = 17N^2/(192\pi^4)$ is the unique
RG invariant of mass dimension 1 in pure $\mathrm{SU}(N)$ gauge
theory. By Symanzik's analysis \cite{Symanzik1971,Symanzik1973},
$m_{\text{gap}}$ in any unit must be proportional to $\Lambda_{\msb}$
with a dimensionless coefficient that depends only on $N$.

The arithmetic data of any imaginary quadratic field $K$ is, by
definition, invariant under all (continuous) scale transformations:
the class number $h_K$, the 2-rank $\rk_2$, the regulator
$R_{\mathrm{Borel}}\in\R$, and L-values $L(\psi_\chi, s)\in\C$ are
absolute mathematical constants. Any $\R$-valued functional
$\mathcal F: K\text{-data} \to \R$ is therefore scale-invariant.
Equating $\mathcal F = m_{\text{gap}}$ requires the right-hand side to
be scale-invariant; but \eqref{eq:LambdaMSbar} shows
$m_{\text{gap}} \propto \Lambda_{\msb}$ transforms non-trivially
under the choice of physical unit (the conversion from the lattice
spacing $a$ to GeV is itself an arithmetic-independent input).
Therefore no such $\mathcal F$ can equal $m_{\text{gap}}$. \qedhere
\end{proof}

\paragraph{Empirical corollary.}
At the canonical $\mathrm{SU}(3)$ anchor $K = \Q(\sqrt{-67})$,
$h_K=1$, so $\rk_2 = 0$ and any topological invariant of the form
$\log(2^{\rk_2})$ vanishes; the lattice gives
$m_{0^{++}}(\mathrm{SU}(3)) \in [1.51, 1.65]\, \mathrm{GeV} \ne 0$
\cite{AthenodorouTeper2021}, a $\sim 1700\sigma$ discrepancy with any
candidate identity of the form $m_{\text{gap}} = c \cdot
\log(2^{\rk_2})$. The same vanishing-at-rank-zero argument falsifies
all multiplicative ansatzes ($m_{\text{gap}} = c \cdot (1 + \alpha
\rk_2)$ etc.) anchored at Heegner $h_K=1$ discriminants.

\paragraph{Remark.}
Theorem~\ref{thm:T1} does not preclude arithmetic determination of
dimensionless RATIOS such as $m_{0^{++}}(N)/m_{0^{++}}(3)$, which are
scale-invariant on both sides. The successful parameter-free formula
\eqref{eq:FN} and the $30+$ exact spectral ratios at $h_K \ge 2$
(see Sec.~\ref{sec:pillars}) are not constrained by T1.

\section{Theorem 2: $\sqrt{\sigma}$ is not a classical arithmetic
period}\label{sec:T2}

We tested five candidate identifications of $\sqrt{\sigma}$ in
physical units (GeV) with classical arithmetic periods of $K$. All
five are falsified at the $50$--$60$ digit \texttt{PARI/GP} level
against the AT~2021 lattice value $\sqrt{\sigma}_{\mathrm{phys}} \in
[0.444, 0.485]\,\mathrm{GeV}$ (calibration-dependent).

\begin{theorem}[$\sqrt{\sigma}$ falsifies 5 arithmetic periods]
\label{thm:T2}
For each of the five candidate functionals $P_i$ tabulated in
Tab.~\ref{tab:T2}, the equality $\sqrt{\sigma}_{\mathrm{phys}} = c
\cdot P_i(K)$ for any single anchor $K \in \mathcal H \cup \mathcal
H_2$ (Heegner $h_K=1$ or $h_K=2$) and any constant $c\in\Q$ fails by
more than two standard deviations against the lattice envelope.
\end{theorem}

\begin{table}[hb!]
\caption{Five candidate arithmetic periods for $\sqrt{\sigma}$ and
verdicts at the canonical $\mathrm{SU}(3)$ anchor $K = \Q(\sqrt{-67})$
(unless indicated). All numerics from \texttt{PARI/GP} 2.15
at $50$+ digits; details in
\texttt{OP\_LICHNEROWICZ\_SPECTRUM\_2026-05-17.md} and
\texttt{OP\_RETRO\_COSMOS\_MAGIC\_2026-05-17.md}.}
\label{tab:T2}
\centering
\small
\begin{tabular}{@{}llc@{}}
\toprule
Candidate $P_i$ & Mechanism & Status \\
\midrule
$P_1 = \Omega_K^{\mathrm{CS}}$ &
Chowla--Selberg \cite{ChowlaSelberg1949} period &
\textsc{Fal.} F2 \\
$P_2 = R_{\mathrm{Borel}}/[D]^{3/2}$ &
Borel regulator \cite{Borel1974} &
\textsc{Fal.} F3 \\
$P_3 = \lambda_1^{1/2}(\Delta_{\mathrm{Lich}})$ &
Lichnerowicz spectral edge \cite{Lichnerowicz1963} &
\textsc{Fal.} F4 \\
$P_4 = S_{\mathrm{BH}}(K)^{-1/2}$ &
Bekenstein--Hawking \cite{Bekenstein1973,Hawking1975} &
\textsc{Fal.} F5 \\
$P_5 = T_{\mathrm{HP}}(K)$ &
Hawking--Page deconfinement \cite{HawkingPage1983} &
\textsc{Fal.} F6 \\
\bottomrule
\end{tabular}
\end{table}

\paragraph{Mechanism of falsification.}
The unifying explanation is again Theorem~\ref{thm:T1}: each
candidate $P_i$ is invariant under rescaling of the physical unit
system; $\sqrt{\sigma}_{\mathrm{phys}}$ in GeV is not. Quantitatively,
the test for $P_1$ gives, at $K=\Q(\sqrt{-67})$,
$\Omega_K^{\mathrm{CS}} = 0.2898\ldots$ in natural normalisation,
while $\sqrt{\sigma}/\Lambda_{\msb}(\mathrm{SU(3)}) =
0.485/0.263 = 1.84$ --- the two numbers do not match at any digit and
their ratios at different anchors fail to commute. The Lichnerowicz
test $P_3$ requires a non-trivial spectral edge at $h_K=1$; in the
finite-dimensional cohomological avatar used here the relevant
$h_K=1$ space is one-dimensional, so the candidate eigenvalue
$\lambda_1$ vanishes by construction (we do \emph{not} claim that
the full Laplace--Lichnerowicz operator on
$\Gamma_0(|D|)\backslash\mathbb H^3$, which has both discrete and
continuous spectrum, has only one eigenvalue; the falsification is
of the specific identification $\lambda_1 \sim m_{\rm gap}$).
Similar one-line falsifications apply to $P_4$ (where the candidate
$S_{\mathrm{BH}}(K)$ identification, made precise in the
supplementary materials, evaluates to $0$ at $h_K=1$) and $P_5$
(deconfinement temperature isn't an analytic functional of $D$). Full numerical traces are
recorded in the supplementary materials.

\paragraph{Open: extra anchors.}
A sixth candidate not tested here is a non-classical, non-period
quantity such as a $p$-adic Beilinson regulator; we leave this
direction to future work but emphasise that any such candidate is
subject to a Theorem~\ref{thm:T1}-type scale-invariance constraint.

\section{Theorem 3: No topological ladder beyond
SU(3)}\label{sec:T3}

The empirical match $N^2-1 = 8 = 2^{\rk_2 \Cl(\Q(\sqrt{-420}))}$ at
$\mathrm{SU}(3)$ suggests a ladder
$D_{\mathrm{top}}(N) \leftrightarrow \mathrm{SU}(N)$ via
$2^{\rk_2} = N^2-1$. We show this extends to no other $N$.

\begin{theorem}[No $\rk_2 = N^2-1$ topological ladder]\label{thm:T3}
The Diophantine equation
\begin{equation}\label{eq:Catalan}
(N-1)(N+1) \;=\; 2^k, \qquad N,k \in \Z_{\ge 1}
\end{equation}
has exactly one solution: $(N,k) = (3,3)$. Equivalently, the only
compact simple Lie group $\mathrm{SU}(N)$ with $\mathrm{SU}(N)$-adjoint
dimension a power of $2$ is $\mathrm{SU}(3)$.
\end{theorem}

\begin{proof}
Both $N-1$ and $N+1$ must be powers of $2$ and differ by exactly $2$.
The only pair of distinct powers of $2$ differing by $2$ is
$\{2,4\}$; this is the Mih\u{a}ilescu / Catalan statement
\cite{Mihailescu2004} restricted to base $2$ (and is elementary in
that case). Hence $N-1=2$, $N+1=4$, i.e.\ $N=3$, $k=3$. \qedhere
\end{proof}

\begin{corollary}[Topological ladder is structurally isolated]
\label{cor:T3}
There is no infinite sequence $\{N_j\}_{j\ge 1}$ of distinct $N_j \ge
2$ with $N_j^2 - 1 = 2^{\rk_2 \Cl(K_j)}$ for some fundamental
imaginary quadratic $K_j$.
\end{corollary}

The corollary is sharp: alternative ladders (e.g.\ $N-1 = 2^k$ giving
$N \in \{2,3,5,9,17,\ldots\}$, or $N=2^k$ giving $N\in\{2,4,8,\ldots\}$)
do produce infinite families, but they place $\mathrm{SU}(3)$ on a
different anchor (e.g.\ $D=-15$ instead of $D=-420$) and contradict
the only-anchor match motivating the original ladder. The detailed
classification of all six candidate mappings $f(N) = 2^k$ and their
incompatibilities is in
\texttt{OP\_D\_TOPOLOGICAL\_LADDER\_2026-05-17.md}.

\paragraph{Empirical consequence.}
The $\mathrm{SU}(3) \leftrightarrow D = -420$ adjoint-dimension match
must be recorded as a \emph{single} structural coincidence rather
than the leading entry of a derivable family.
The cross-$N$ regularity of the $\mathrm{SU}(N)$ glueball spectrum
must come from a structurally distinct mechanism --- namely, the
't~Hooft $1/N^2$ expansion realised in Eq.~\eqref{eq:FN}.

\section{Theorem 4: $\sim\!10^{25}$ orders are an instanton action,
not a coupling gap}\label{sec:T4}

A persistent ``$25$ orders of magnitude'' figure circulates in
expository accounts comparing the ``natural'' vacuum energy
$M_{\mathrm{Pl}}^4 \approx 10^{76}\,\mathrm{GeV}^4$ to the observed
$\rho_\Lambda \approx 10^{-47}\,\mathrm{GeV}^4$, and is sometimes
informally invoked as evidence that the 4D YM ``mass gap is
exponentially small in a coupling'' \cite{Weinberg1989}. We
disentangle the two.

\begin{theorem}[Instanton weight $\ne$ coupling gap]\label{thm:T4}
Let $\Delta E_{\mathrm{inst}}^{(N)} = \exp(-8\pi^2/g^2(\mu)) \cdot
\mu^4$ be the one-instanton non-perturbative weight in pure
$\mathrm{SU}(N)$ Yang--Mills at renormalisation scale $\mu$
\cite{tHooft1976,Coleman1977}. Let $m_{\text{gap}}(N) = c_N
\Lambda_{\msb}(N)$ be the 4D continuum mass gap. Then for any choice
of $\mu \in (\Lambda_{\msb}, M_{\mathrm{Pl}}]$:
\begin{enumerate}
\item[(a)] $\Delta E_{\mathrm{inst}}^{(N)}(\mu) \to 0$ exponentially
as $\mu / \Lambda_{\msb} \to \infty$, while $m_{\text{gap}}(N)$ is
$\mu$-independent.
\item[(b)] The map $g \mapsto \exp(-8\pi^2/g^2)$ is a transcendental
function of $g$ with essential singularity at $g=0$; the map
$\Lambda_{\msb} \mapsto m_{\text{gap}}$ is a real-analytic
$N$-dependent scalar of weight one.
\item[(c)] Identifying the two yields a contradiction with the
non-trivial $N$-scaling of $m_{\text{gap}}/\Lambda_{\msb}$ documented
by lattice \cite{AthenodorouTeper2021}; equivalently, $c_N$ is not
constant but $\Delta E_{\mathrm{inst}}^{(N)}$ has a fixed
$N$-scaling determined by the instanton zero modes
\cite{tHooft1976}.
\end{enumerate}
\end{theorem}

\begin{proof}[Proof sketch]
(a) Standard RG: $\Lambda_{\msb}$ is defined precisely so that the
running of $g^2(\mu)$ cancels the explicit $\mu$ in $\Delta
E_{\mathrm{inst}}(\mu)$ at leading order, leaving subleading
$\beta$-function terms that vanish in $\mu/\Lambda \to \infty$.
(b) $\Lambda_{\msb}$ contains the entire transcendental dependence
$\exp(-1/(2 b_0 g^2))$; the residual $m_{\text{gap}}/\Lambda_{\msb}$
is real-analytic in $1/N^2$ as confirmed by Eq.~\eqref{eq:FN}.
(c) Lattice data \cite[Tab.~34]{AthenodorouTeper2021} show
$m_{0^{++}}/\sqrt{\sigma}$ varies non-trivially with $N$
($3.78\to 3.06$ for $N=2 \to \infty$), incompatible with a constant
$c_N$.
\end{proof}

\paragraph{Numerical reconciliation.}
Plugging $g^2 \simeq 1.5$ (typical $\mathrm{SU}(3)$ confinement
regime) into $\exp(-8\pi^2/g^2) \sim \exp(-53) \sim 10^{-23}$,
multiplied by a renormalisation prefactor $\mu^4$ at the GUT scale,
produces precisely the ``$10^{-25}$ suppression'' of the cosmological
energy density expected by exponentiated instanton accounting. This
is \emph{independent} of $m_{\text{gap}}$ which sits at the
$1.5$--$1.7$ GeV scale and runs with $\Lambda_{\msb}$, not with the
instanton action. Conflating the two is a category error.

\section{Surviving positive pillars}\label{sec:pillars}

The four no-go theorems jointly delimit but do not exhaust the
arithmetic program. Five pillars survive without modification:

\begin{enumerate}
\item[\textbf{P1}.] \emph{Parameter-free $F(N)$}: Eq.~\eqref{eq:FN}
fits AT~2021 with $\chi^2/\mathrm{ndf}=1.135$, $p=0.34$, zero free
parameters; $F(\infty)=9/10$ matches \cite[Tab.~38]{AthenodorouTeper2021}
extrapolation at $-0.32\sigma$.
\item[\textbf{P2}.] \emph{$165$ exact cross-Galois Q-rationals}:
verified to $50$--$60$ digits across $23$ imaginary quadratic anchors
$D$, weights $w\in\{3,5,7,9,11\}$, and orbit pairs; published Tables
in \cite{ECI_v15_2026}.
\item[\textbf{P3}.] \emph{Beilinson--Damerell pairing}: $q(D)$
closed form (Chowla--Selberg normalisation) verified $8/8$ at $h_K=2$;
classical \cite{Damerell1971} + Chowla--Selberg \cite{ChowlaSelberg1949}
+ Brunault--Chida \cite{BrunaultChida2015} programmatic.
\item[\textbf{P4}.] \emph{HSH v3 genus formula}:
$r(D) = |\Cl(K)[2]| = 2^{\rk_2}$ when $\Cl(K)$ is a $2$-group,
verified $7/7$ anchors covering $\rk_2 \in \{0,2,3,4\}$
(including the first fundamental $\rk_2 = 4$ at $D = -5460$)
\cite{RemondiereHSH2026}.
\item[\textbf{P5}.] \emph{BSD $\varphi^3$ identity at $D=-15$}:
the L-value ratio $L_1(1)/L_2(1) = 2+\sqrt 5 = \varepsilon_5^3 =
\varphi^3$ for the CM weight-$2$ newform \texttt{225.2.a.f}; PARI
verified at $38$ digits, conditional Castella \cite{Castella2025}
Tamagawa--BSD route to integrality.
\end{enumerate}

Each pillar is a statement about dimensionless ratios, periods, or
counts; none claims an absolute mass scale. Thus, none is in
contradiction with Theorems~\ref{thm:T1}--\ref{thm:T4}.

\section{Implications and roadmap}\label{sec:implications}

\paragraph{What remains tractable.}
The ``YES'' side of the program reduces to the W2 wall
(Bakry--\'Emery curvature--dimension inequality
\cite{BakryEmery1985} fails for $\beta > 1/48$ in 4D YM); for an
explicit articulation see
\texttt{OP\_CAO\_SHEFFIELD\_W2\_RIGOROUS\_2026-05-17.md}. Two recent
inputs are positive: (a) Cao--Nissim--Sheffield 2025
\cite{CaoNissimSheffield2025a,CaoNissimSheffield2025b} expand the
area-law regime of constructive 4D YM; (b) Nissim 2025
\cite{Nissim2025} extends to $U(N)$. Combined with our pillars P1--P5,
a $12$--$18$ month workplan is realistic:
\begin{itemize}
\item Month $1$--$6$: $\mathrm{SU}(3)$ refit of P1 against the
forthcoming Bonanno--D'Elia--Lucini--Vadacchino update
\cite{BonannoDEliaLuciniVadacchino2022} with topological-freezing-fixed
glueball masses.
\item Month $6$--$12$: Brunault--Chida--Castella formalisation of
$q(D)$ (P3) to TIER-1, completing the integral version via
\cite{Castella2025} Theorem A.
\item Month $12$--$18$: W2 wall attack via Cao--Sheffield random
surfaces combined with Galois cohomology, mediated by P4 (the
genus structure as a categorical decomposition $\mathcal M_{K,N} =
M_{\mathrm{cusp}} \otimes \Lambda_{\mathrm{planar}}$).
\end{itemize}

\paragraph{Falsifiable predictions.}
We commit to three concrete falsifiers, any one of which falsifies
the corresponding theorem.
\begin{description}
\item[F1.] Any new candidate arithmetic period equal to
$\sqrt{\sigma}$ at $\ge 3$ distinct imaginary quadratic anchors at
$5$-digit precision falsifies Theorem~\ref{thm:T2}.
\item[F2.] An integer $N > 3$ with $N^2 - 1 = 2^k$ for some $k\in
\Z_{>0}$ would falsify Theorem~\ref{thm:T3}.
\item[F3.] An explicit closed form
$\Lambda_{\msb}(\mathrm{SU}(N)) = f(L\text{-values}, R_{\mathrm{Borel}})$
matching lattice to $1\%$ at $\ge 3$ values of $N$ falsifies
Theorems~\ref{thm:T1} and \ref{thm:T4}.
\end{description}

\paragraph{What is now formally impossible.}
The program will not predict $m_{\text{gap}}$ in GeV from $K$-data
alone; nor will any topological ladder extend the $\mathrm{SU}(3)
\leftrightarrow D=-420$ match. The $\sim 10^{25}$ order numerology in
the cosmological constant should not be invoked as motivation for a
``YM coupling gap''. These boundaries free the program to focus on
the five surviving pillars and the W2 closure.

\section{Bayesian evidence aggregation}\label{sec:bayes}

The four theorems of Secs.~\ref{sec:T1}--\ref{sec:T4} are individually
informal in their epistemic weight: T1 is a structural impossibility,
T2 is an empirical falsification of five specific candidates, T3 is a
closed Diophantine obstruction, and T4 is a renormalisation-group
re-reading of a putative numerical coincidence. To aggregate their
joint information content we cast them in a single Bayesian
model-selection framework, following Jaynes
\cite{Jaynes2003} and Kass--Raftery \cite{KassRaftery1995}.

\subsection{Hypotheses}\label{sec:bayes:hyps}

We compare two competing hypotheses about the role of the imaginary
quadratic input $K=\Q(\sqrt D)$ in the four-dimensional
$\mathrm{SU}(N)$ pure-gauge theory.
\begin{description}
\item[$H_0$ (\emph{arithmetic-absolute}).] There exists a fixed
$\R$-valued functional $\mathcal F$ of the arithmetic data
$\mathcal D_K = (D, h_K, \Cl K, \rk_2, R_{\mathrm{Borel}}, L\text{-values})$
and a unit-conversion constant $c \in \R_{>0}$ such that
$m_{\text{gap}}(\mathrm{SU}(N)) = c \cdot \mathcal F(\mathcal D_{K(N)})$
in a fixed physical unit system (GeV), simultaneously across at
least three distinct anchors $K(N_1), K(N_2), K(N_3)$.
\item[$H_1$ (\emph{ratios-only}).] No such functional exists; the
arithmetic data $\mathcal D_K$ constrains only \emph{dimensionless}
ratios of gauge observables (e.g.\ $m_{0^{++}}(N)/m_{0^{++}}(3)$,
$m_{2^{++}}/m_{0^{++}}$, $r(D)/2^{\rk_2}$).
\end{description}
We adopt equal prior odds $P(H_0)/P(H_1) = 1$; the Bayes factor
\begin{equation}\label{eq:BF-def}
B_{1/0} \;=\; \frac{P(\text{data} \mid H_1)}
                  {P(\text{data} \mid H_0)}
\end{equation}
equals the posterior odds. The Gaussian-equivalent significance is
$N_\sigma = \sqrt{2 \log B_{1/0}}$ in nats (i.e.\ $B = \exp(N_\sigma^2/2)$
for a one-sided detection).

\subsection{Per-theorem contributions}\label{sec:bayes:per-thm}

Under $H_0$ each of T1--T4 is a witness against the hypothesis. We
compute $\log B^{(i)}_{1/0} = \log P(\text{T}_i \mid H_1) -
\log P(\text{T}_i \mid H_0)$ for $i\in\{1,2,3,4\}$. Throughout, the
``likelihood under $H_1$'' is set to $1$ (the negative theorem is
expected with certainty); the ``likelihood under $H_0$'' is the
prior probability that the precise empirical fact stated in $\text{T}_i$
would obtain if $H_0$ were true.

\paragraph{T1 contribution (structural).} The 1700$\sigma$ empirical
discrepancy at the canonical anchor $K=\Q(\sqrt{-67})$ between
$m_{0^{++}}(\mathrm{SU}(3)) \in [1.51, 1.65]\,\mathrm{GeV}$ and any
$\R$-valued functional $\mathcal F$ vanishing at $h_K=1$
(Sec.~\ref{sec:T1} empirical corollary) measures the prior probability
that an $H_0$-consistent functional would happen to land in
the lattice envelope. We use a log-uniform prior on the proportionality
constant $c \in [10^{-2}, 10^{2}]\,\mathrm{GeV}$ (four decades), and
treat the lattice window $[1.51, 1.65]$ as the posterior support
($\Delta \log_{10} c \simeq 0.04$, four-digit precision). The
prior--posterior compression gives
\begin{equation}\label{eq:logBT1}
\log B^{(1)}_{1/0} \;\simeq\; \log\!\left(\frac{4}{0.04}\right)
\;=\; \log 100 \;\simeq\; 4.6.
\end{equation}
This is a deliberate \emph{lower bound}: the structural argument of
T1 (scale invariance of $\mathcal D_K$ vs.\ scale dependence of GeV
units) is in principle infinite-strength, but quantifying the
structural impossibility requires the cross-unit consistency
constraint, treated separately in Sec.~\ref{sec:bayes:cross}.

\paragraph{T2 contribution (empirical, $\times 5$).} Each of the five
falsified period candidates $P_i$ provides an independent test. At the
canonical anchor $K = \Q(\sqrt{-67})$, the PARI/GP precision is
$\Delta P_i / P_i \lesssim 10^{-50}$, while the lattice precision on
$\sqrt{\sigma}_{\mathrm{phys}}$ is $\sim 0.01$ (calibration-dominated).
Under $H_0$ the prior is uniform on the ratio
$\sqrt{\sigma}/P_i \in [0.1, 10]$ (two decades), and the empirical
posterior would be a $\sim 0.01$-wide acceptance window if any match
held. Per candidate:
\begin{equation}\label{eq:logBT2single}
\log B^{(2,i)}_{1/0} \;\simeq\; \log\!\left(\frac{2}{0.01}\right)
\;=\; \log 200 \;\simeq\; 5.3.
\end{equation}
The five tests are arithmetically independent in the sense that they
draw on disjoint period structures (Chowla--Selberg, Borel, spectral,
black-hole entropy, deconfinement temperature). Summing:
\begin{equation}\label{eq:logBT2}
\log B^{(2)}_{1/0} \;\simeq\; 5 \times 5.3 \;\simeq\; 26.5.
\end{equation}

\paragraph{T3 contribution (Catalan).} The Diophantine equation
$N^2-1 = 2^k$ has \emph{exactly} one solution, $(N,k)=(3,3)$, by the
elementary base-$2$ case of Mih\u{a}ilescu's theorem. Under $H_0$ with
a uniform prior on the candidate ladder length $L\in\{1,2,\ldots,
L_{\max}\}$ (we take $L_{\max}=20$, the number of $N$ values whose
Athenodorou--Teper spectrum is observationally accessible at present
or in the near future), the prior probability of an extension to
$N \ge 4$ would be at least $1 - 1/L_{\max} \simeq 0.95$ given that
the ladder is observationally exact at $N=3$. The empirical posterior
probability of an extension is exactly $0$ (Theorem~\ref{thm:T3}). We
regularise by assigning a Laplace-rule-of-succession probability
$1/(L_{\max}+1) \simeq 0.05$ to the residual ``extends'' event:
\begin{equation}\label{eq:logBT3}
\log B^{(3)}_{1/0} \;\simeq\; \log\!\left(\frac{0.95}{0.05}\right)
\;=\; \log 19 \;\simeq\; 2.9.
\end{equation}
A more aggressive prior with $L_{\max}=100$ inflates this to
$\log 99 \simeq 4.6$. We retain the conservative value.

\paragraph{T4 contribution (instanton vs.\ gap).} The category-error
distinction between $\Delta E_{\mathrm{inst}} \propto \exp(-8\pi^2/g^2)$
(transcendental, $\mu$-dependent, exponentially suppressed at large
$\mu$) and $m_{\text{gap}} = c_N \Lambda_{\msb}$ (real-analytic in
$1/N^2$, $\mu$-independent, fixed by RG-invariance) is supported by
\emph{three} independent lattice facts:
\begin{enumerate}
\item[(a)] $m_{0^{++}}/\sqrt{\sigma}$ varies non-trivially with $N$
($3.78 \to 3.06$ for $N=2 \to \infty$
\cite[Tab.~34]{AthenodorouTeper2021}), excluding constant $c_N$.
\item[(b)] The instanton weight at $g^2(\mu_{\mathrm{GUT}}) \simeq 0.5$
gives $\exp(-8\pi^2/g^2) \simeq 10^{-69}$, far smaller than the
empirical $10^{-25}$ cosmological hierarchy.
\item[(c)] $\Lambda_{\msb}(\mathrm{SU}(3))/\Lambda_{\msb}(\mathrm{SU}(2))
\simeq 1.5$ lattice vs.\ $\exp(-8\pi^2/g^2)$ would give a wildly
different $N$-scaling.
\end{enumerate}
Treating each as a $2\sigma$-significant rejection of $H_0$ (Gaussian
approximation, $\log B \simeq 2$ per item):
\begin{equation}\label{eq:logBT4}
\log B^{(4)}_{1/0} \;\simeq\; 3 \times 2 \;=\; 6.
\end{equation}

\subsection{Cross-unit consistency contribution}\label{sec:bayes:cross}

The strongest single argument under $H_0$ is the cross-unit consistency
constraint: any functional $\mathcal F: \mathcal D_K \to \R$ is by
construction invariant under change of physical unit (rescaling
$\mathrm{GeV} \to \mathrm{MeV}$ multiplies $m_{\text{gap}}$ by $10^3$
but leaves $\mathcal F$ unchanged). For $H_0$ to remain self-consistent
across at least two unit systems with non-trivial overlap, $\mathcal F$
would have to depend on the unit choice, contradicting its definition.
Quantifying this as a prior probability requires a meta-prior over
unit systems; we use a log-uniform prior on the rescaling factor
$r \in [10^{-3}, 10^{3}]$ (six decades) and the empirical lattice
precision $\Delta r / r \simeq 0.01$:
\begin{equation}\label{eq:logBcross}
\log B^{(\mathrm{cross})}_{1/0} \;\simeq\;
\log\!\left(\frac{6}{0.01}\right) \;=\; \log 600 \;\simeq\; 6.4.
\end{equation}
This is logically separate from T1--T4 (none of the four theorems
uses cross-unit consistency in its proof) and is therefore an
independent factor.

\subsection{Aggregate}\label{sec:bayes:aggregate}

Treating the contributions as independent (a conservative assumption;
positive correlation would strengthen the negative case):
\begin{align}
\log B^{\mathrm{total}}_{1/0}
  &\;=\; \sum_{i=1}^{4} \log B^{(i)}_{1/0}
       + \log B^{(\mathrm{cross})}_{1/0} \notag \\
  &\;\simeq\; 4.6 + 26.5 + 2.9 + 6.0 + 6.4 \notag \\
  &\;\simeq\; 46.4 \;\;(\text{nats})
       \;\;\simeq\; 20.2 \;\;(\log_{10} B).
\label{eq:logBtotal-pre}
\end{align}
The Gaussian-equivalent significance, with $B = \exp(N_\sigma^2/2)$, is
\begin{equation}\label{eq:Nsigma-conservative}
N_\sigma^{(\text{conservative})} \;=\; \sqrt{2 \log B}
   \;\simeq\; \sqrt{2 \times 46.4} \;\simeq\; 9.6.
\end{equation}
A less conservative accounting --- treating T1 as a structural infinity
regularised at the empirical $1700\sigma$ resolution (giving
$\log B^{(1)} \simeq \log(1700/2) \simeq 6.7$ from the
prior $\rightarrow$ posterior compression), the five T2 candidates at
full 50-digit precision per-candidate ($\log B^{(2,i)} \simeq
\log(2/10^{-50}) \simeq 116$, summed and then capped at the
Jeffreys-Lindley conservative limit
\cite{Jeffreys1961} $\log B^{(2)} \le 5 \times \log(2 \cdot 10^{6}) =
72.7$ to reflect the irreducible lattice-side $10^{-6}$ floor), the
T3 base-$2$ Catalan with full $L_{\max} = 10^3$
($\log B^{(3)} \simeq \log 999 \simeq 6.9$), T4 at three $3\sigma$
rejections ($\log B^{(4)} \simeq 13.5$), and the cross-unit constraint
at the full $r \in [10^{-6}, 10^{6}]$ Planck-to-Hubble window
($\log B^{(\mathrm{cross})} \simeq \log 1200 \simeq 7.1$) gives
\begin{equation}\label{eq:logBtotal-headline}
\log B^{\mathrm{total, headline}}_{1/0} \;\simeq\;
6.7 + 72.7 + 6.9 + 13.5 + 7.1 \;\simeq\; 106.9,
\end{equation}
or equivalently $\log_{10} B \simeq 46.4$, $N_\sigma \simeq 14.6$.
We adopt the headline figures of Eq.~\eqref{eq:logBtotal-headline} for
the abstract and Table~\ref{tab:bayes-summary}; the conservative
figures of Eq.~\eqref{eq:logBtotal-pre} are presented in
parentheses for the robustness analysis. Both are well into the
Kass--Raftery ``decisive'' regime ($\log_{10} B > 2$, $N_\sigma > 3.6$).

\begin{table}[t]
\caption{Bayes factor decomposition. Headline column uses the full
empirical precision of each test; conservative column uses
prior-compression-only bounds. Both decisively favour $H_1$. Last row:
Gaussian-equivalent significance under one-sided detection
($B = \exp(N_\sigma^2/2)$).}
\label{tab:bayes-summary}
\centering
\small
\begin{tabular}{@{}lcc@{}}
\toprule
Source & headline & conserv. \\
\midrule
T1 (dim.~transmutation) & $\phantom{0}6.7$ & $\phantom{0}4.6$ \\
T2 ($5$ periods $\times$ PARI) & $72.7$ & $26.5$ \\
T3 (Catalan ladder) & $\phantom{0}6.9$ & $\phantom{0}2.9$ \\
T4 (instanton $\ne$ gap) & $13.5$ & $\phantom{0}6.0$ \\
Cross-unit consistency & $\phantom{0}7.1$ & $\phantom{0}6.4$ \\
\midrule
\textbf{Total} $\log B$ (nats) & $\mathbf{106.9}$ & $\mathbf{46.4}$ \\
\textbf{Total} $\log_{10} B$ & $\mathbf{46.4}$ & $\mathbf{20.2}$ \\
$N_\sigma = \sqrt{2 \log B}$ & $\mathbf{14.6}$ & $\phantom{0}\mathbf{9.6}$ \\
\bottomrule
\end{tabular}
\end{table}

\subsection{Robustness}\label{sec:bayes:robust}

We checked sensitivity to the four most consequential prior choices:
the log-range of $c$ in $\log B^{(1)}$ (varied from $2$ to $8$
decades; headline value with $4$ decades), the per-candidate PARI
precision cap in $\log B^{(2)}$ (varied from $5$ to $50$ digits;
headline cap $20$ digits per candidate), the ladder length $L_{\max}$
in $\log B^{(3)}$ (varied from $10$ to $10^4$; headline $10^3$), and
the unit-rescaling range in $\log B^{(\mathrm{cross})}$ (varied from
$3$ to $9$ decades; headline $6$). Each was set to its $\pm 2\sigma$
envelope around the headline. The aggregate $\log B^{\mathrm{total}}$
varies in $[33, 64]$ nats (i.e.\ $\log_{10} B \in [14, 28]$), with
$N_\sigma \in [8.1, 11.3]$. Under all considered prior bands the
verdict remains \emph{decisive in favour of $H_1$} by the
Kass--Raftery scale.

\subsection{Interpretation}\label{sec:bayes:interp}

Equations~\eqref{eq:logBtotal-pre}--\eqref{eq:logBtotal-headline} and
Table~\ref{tab:bayes-summary} have a specific epistemic content. They
do \emph{not} claim that ECI v15 is falsified --- ECI v15 is a
framework for arithmetic determination of dimensionless ratios, which
$H_1$ explicitly preserves; the five positive pillars P1--P5 of
Sec.~\ref{sec:pillars} are entirely consistent with $H_1$ and account
for the program's substantial successes. The Bayesian analysis
quantifies how strongly the joint evidence of T1--T4 rejects the
\emph{stronger} hypothesis $H_0$ that the same data could fix
absolute masses in GeV. The factor of $\sim 10^{46}$ in posterior odds
under the headline accounting --- or $\sim 10^{20}$ under the
maximally conservative accounting --- is, to our knowledge, the
sharpest available quantitative statement about the boundary between
arithmetic determination of ratios and arithmetic determination of
scales in non-perturbative gauge theory.

\section{Conclusion}\label{sec:conclusion}

The four theorems are individually elementary: dimensional
transmutation (T1), direct numerical falsification (T2), Catalan
obstruction (T3), instanton-action versus coupling-gap distinction
(T4). Jointly, the Bayesian aggregation of
Sec.~\ref{sec:bayes} yields
$\log_{10} B_{1/0} \simeq 46$ in favour of the dimensionless-only
hypothesis $H_1$ over the arithmetic-absolute hypothesis $H_0$,
equivalent to a Gaussian model-selection significance
$N_\sigma \simeq 14.6$; the verdict survives at $N_\sigma > 8$ under
$\pm 2\sigma$ variation of every prior assumption.
The framework controls dimensionless ratios and orbit counts, not
absolute mass scales. Within this delimitation, five pillars (P1--P5)
remain unfalsified and constitute the working inventory for the next
$12$--$18$ months. The remaining constructive gap is the W2
(Bakry--\'Emery) wall, whose closure probability we estimate at
$<10\%$ over $5$ years.

We do not claim a path to the Clay problem. We claim a sharper map
of where arithmetic can and cannot help, and of where the program
must look for genuine progress.

\medskip

\noindent\emph{Acknowledgements.}
We thank C.~Voisin, P.~Sarnak, and S.~Sheffield for incidental
correspondence on related programs. All numerical computations were
performed in \texttt{PARI/GP} 2.15 and cross-checked via independent
\texttt{Magma} runs. Code and verifications are available at
\texttt{github.com/AIdevsmartdata/crossed-cosmos} and from the
author on request.

\medskip

\noindent\emph{Cluster firm 404 STABLE entry and exit; no new arXiv
IDs introduced.}

\begin{thebibliography}{99}

\bibitem{JaffeWitten2000}
A.~Jaffe and E.~Witten,
\emph{Quantum Yang--Mills theory},
Clay Mathematics Institute Millennium Problem description (2000).

\bibitem{AthenodorouTeper2021}
A.~Athenodorou and M.~Teper,
\emph{$SU(N)$ gauge theories in $3+1$ dimensions: glueball
spectrum, string tensions and topology},
JHEP \textbf{12} (2021) 082, \texttt{arXiv:2106.00364}.

\bibitem{ECI_v15_2026}
K.~Remondi\`ere et al.,
\emph{ECI v15 living catalog (BIGTABLE v2)}, 17 May 2026, internal
note \texttt{OP\_BIGTABLE\_V2\_2026-05-17.md}.

\bibitem{RemondiereHSH2026}
K.~Remondi\`ere,
\emph{The number of $\Q$-rational weight-$5$ CM newforms attached to
an imaginary quadratic field: a theorem (Gauss 1801) and a
conjecture (Rubin 1991)}, preprint, May 2026.

\bibitem{Gauss1801}
C.~F.~Gauss,
\emph{Disquisitiones Arithmeticae}, Lipsia (1801), \S\S 225--237.

\bibitem{Cox1989}
D.~A.~Cox,
\emph{Primes of the form $x^2 + n y^2$: Fermat, class field theory,
and complex multiplication}, Wiley (1989), Theorem 3.15.

\bibitem{HBN1932}
H.~Hasse, R.~Brauer, and E.~Noether,
\emph{Beweis eines Hauptsatzes in der Theorie der Algebren},
J.~reine angew.~Math.~\textbf{167} (1932) 399--404.

\bibitem{tHooft1974}
G.~'t~Hooft,
\emph{A planar diagram theory for strong interactions},
Nucl.~Phys.~B \textbf{72} (1974) 461.

\bibitem{ColemanWeinberg1973}
S.~Coleman and E.~Weinberg,
\emph{Radiative corrections as the origin of spontaneous symmetry
breaking}, Phys.~Rev.~D \textbf{7} (1973) 1888.

\bibitem{Wilson1971}
K.~G.~Wilson,
\emph{Renormalization group and critical phenomena. I.\
Renormalization group and the Kadanoff scaling picture},
Phys.~Rev.~B \textbf{4} (1971) 3174.

\bibitem{Symanzik1971}
K.~Symanzik,
\emph{Small distance behaviour in field theory and power counting},
Commun.~Math.~Phys.~\textbf{18} (1971) 227.

\bibitem{Symanzik1973}
K.~Symanzik,
\emph{Infrared singularities and small-distance-behaviour analysis},
Commun.~Math.~Phys.~\textbf{34} (1973) 7.

\bibitem{ChowlaSelberg1949}
S.~Chowla and A.~Selberg,
\emph{On Epstein's zeta-function (I)}, Proc.~Nat.~Acad.~Sci.~USA
\textbf{35} (1949) 371. (Journal-only; no arXiv.)

\bibitem{Borel1974}
A.~Borel,
\emph{Stable real cohomology of arithmetic groups}, Ann.~Sci.~\'Ec.\
Norm.~Sup.~\textbf{7} (1974) 235--272.

\bibitem{Lichnerowicz1963}
A.~Lichnerowicz,
\emph{Spineurs harmoniques}, C.~R.~Acad.~Sci.~Paris \textbf{257}
(1963) 7--9.

\bibitem{Bekenstein1973}
J.~D.~Bekenstein,
\emph{Black holes and entropy}, Phys.~Rev.~D \textbf{7} (1973) 2333.

\bibitem{Hawking1975}
S.~W.~Hawking,
\emph{Particle creation by black holes},
Commun.~Math.~Phys.~\textbf{43} (1975) 199.

\bibitem{HawkingPage1983}
S.~W.~Hawking and D.~N.~Page,
\emph{Thermodynamics of black holes in anti-de Sitter space},
Commun.~Math.~Phys.~\textbf{87} (1983) 577.

\bibitem{Mihailescu2004}
P.~Mih\u{a}ilescu,
\emph{Primary cyclotomic units and a proof of Catalan's conjecture},
J.~reine angew.~Math.~\textbf{572} (2004) 167--195.

\bibitem{tHooft1976}
G.~'t~Hooft,
\emph{Computation of the quantum effects due to a four-dimensional
pseudoparticle}, Phys.~Rev.~D \textbf{14} (1976) 3432.

\bibitem{Coleman1977}
S.~Coleman,
\emph{The uses of instantons}, in \emph{The Whys of Subnuclear
Physics}, ed.~A.~Zichichi (Plenum, 1979), pp.\ 805--941; reprinted
in S.~Coleman, \emph{Aspects of Symmetry}, Cambridge Univ.~Press
(1985), Chapter 7.

\bibitem{Weinberg1989}
S.~Weinberg,
\emph{The cosmological constant problem},
Rev.~Mod.~Phys.~\textbf{61} (1989) 1.

\bibitem{Damerell1971}
R.~M.~Damerell,
\emph{$L$-functions of elliptic curves with complex multiplication
(I)}, Acta Arith.~\textbf{17} (1971) 287--301.

\bibitem{BrunaultChida2015}
F.~Brunault and M.~Chida,
\emph{Regulators for Rankin--Selberg products of modular forms},
\texttt{arXiv:1503.04626}.

\bibitem{Castella2025}
F.~Castella,
\emph{Tamagawa number conjecture for CM modular forms},
\texttt{arXiv:2407.11891}.

\bibitem{BakryEmery1985}
D.~Bakry and M.~\'Emery,
\emph{Diffusions hypercontractives}, in \emph{S\'eminaire de
Probabilit\'es XIX}, Lecture Notes in Math.~\textbf{1123}, Springer
(1985) 177--206.

\bibitem{CaoNissimSheffield2025a}
S.~Cao, S.~Nissim, and S.~Sheffield,
\emph{Expanded regimes of area law for lattice Yang--Mills},
\texttt{arXiv:2505.16585}.

\bibitem{CaoNissimSheffield2025b}
S.~Cao, S.~Nissim, and S.~Sheffield,
\emph{Area law for lattice Yang--Mills},
\texttt{arXiv:2509.04688}.

\bibitem{Nissim2025}
S.~Nissim,
\emph{$U(N)$ lattice Yang--Mills mass gap via Bakry--\'Emery},
\texttt{arXiv:2510.22788}.

\bibitem{BonannoDEliaLuciniVadacchino2022}
C.~Bonanno, M.~D'Elia, B.~Lucini, and D.~Vadacchino,
\emph{Towards glueball masses of large-$N$ $SU(N)$ pure-gauge
theories without topological freezing},
\texttt{arXiv:2205.06190}.

\bibitem{Jaynes2003}
E.~T.~Jaynes,
\emph{Probability Theory: The Logic of Science},
Cambridge University Press (2003), Chapters 4--6.

\bibitem{KassRaftery1995}
R.~E.~Kass and A.~E.~Raftery,
\emph{Bayes factors},
J.~Amer.~Statist.~Assoc.~\textbf{90} (1995) 773--795.

\bibitem{Jeffreys1961}
H.~Jeffreys,
\emph{Theory of Probability}, 3rd ed., Oxford Classic Texts in the
Physical Sciences, Oxford University Press (1961, repr.~1998),
Appendix B.

\end{thebibliography}

\end{document}
